\section*{From 3-rounds to multiple rounds using deterministic round elimination (author Section 5)}
In this section we show how we can use an oracle for computing $L_{d+1}^B$ in order to compute $L_d^A$. The gadget in our reduction involves a special function, $\XorEq_k$, which is a small modification of the standard Equality function.
\begin{definition}
The function $\XorEq_k:[k]\times([k]\times\{0,1\})\to\{0,1\}$ is given by
\[\XorEq_k(x;y,z)=\begin{cases}z&\text{if }x\ne y,\\1-z&\text{if }x=y.\end{cases}\]
\end{definition}
The key property of $\XorEq_k$ is that it has a fooling set lower bound that is tight. In particular, $\{(x;y,z):x=y\}$ is a fooling set of cardinality $2k$, and there is a 2-round protocol for solving $\XorEq_k$ with $2k$ leaves (where Alice just sends her full input to Bob, and he replies with the output).
\begin{lemma}[Fooling set lower-bound for $\XorEq_k$]\label{lem:fooling}
Let $\pi$ be a protocol for solving the function $f:([k]\times[a])\times([k]\times\{0,1\})\to\{0,1\}$ given by $f(x_0,x_1;y,z)=\XorEq_k(x_0;y,z)$.\footnote{In our definition of $f$, the input $x_1$ does not affect the output of the function. The fact that this lemma holds even when there is an extraneous input like $x_1$ will be used later.} Then
\[\pi(x_0,x_1;y,z)\ne\pi(x'_0,x'_1;y',z')\]
if either
\begin{itemize}
\item $y=x_0\ne x'_0$, or
\item $y=x_0=x'_0$ and $z\ne z'$.
\end{itemize}
\end{lemma}
\begin{proof}
First, suppose for contradiction that $\pi(x_0,x_1;y,z)=\pi(x'_0,x'_1;y',z')$ and $y=x_0\ne x'_0$. Since leaves are combinatorial rectangles, we can infer that $\pi(x'_0,x'_1;y,z)=\pi(x_0,x_1;y,z)$. But since $y=x_0\ne x'_0$, we know that
\[f(x_0,x_1;y,z)=1-z\ne z=f(x'_0,x'_1;y,z)\]
so this contradicts that $\pi(x_0,x_1;y,z)$ is a monochromatic leaf.
Similarly, if $y=x_0=x'_0$ and $z\ne z'$, then we have $f(x_0,x_1;y,z)\ne f(x'_0,x'_1;y',z')$, so $\pi(x_0,x_1;y,z)\ne\pi(x'_0,x'_1;y',z')$ by monochromaticness.
\end{proof}
The main technical portion of our reduction is the following deterministic variant of the round-elimination lemma.
\begin{lemma}[Restatement of Lemma 3]\label{lem:round-elimination}
Let $d\ge3$. Let $f:[a]\times[b]\to\{0,1\}$. Let $F:([k]\times[a])\times([k]\times[b]\times\{0,1\}\times\{0,1\})\to\{0,1\}$ be given by
\[F(x_0,x_1;y_0,y_1,z,i)=\begin{cases}\XorEq_k(x_0;y_0,z)&\text{if }i=0,\\f(x_1;y_1)&\text{if }i=1.\end{cases}\]
Then we have
\[\min\{4k,2k-2+L_d^B(f)\}\le L_{d+1}^A(F)\le2k+L_d^B(f).\]
\end{lemma}
\begin{proof}
The upper bound comes from the protocol where Alice skips the first round of communication, Bob sends $i$ to Alice and begins running the best $d$-round Bob-first protocol for $\XorEq_k$ or $f$, based on whether $i=0$ or $i=1$. In particular, we have that
\[L_{d+1}^A(F)\le L_d^B(\XorEq_k)+L_d^B(f)\le2k+L_d^B(f),\]
where the last upper bound uses that $d\ge3$.
We now argue the lower bound. Suppose $\pi$ is a $(d+1)$-round Alice-first protocol for $F$. Let $\mathcal L=\{\ell_1,\ldots,\ell_{L(\pi)}\}$ denote the set of leaves of $\pi$.
Our arguments split into two cases depending on whether there is a good input $x_0\in[k]$. We say an input $x_0\in[k]$ is good if all of Alice's inputs that begin with $x_0$ are placed in a single partition. More formally, let $\mathcal P=\{P_1,\ldots,P_{|\mathcal P|}\}$ be the partition of Alice's inputs corresponding to the first round of $\pi$. We say $x_0\in[k]$ is good if there exists a $q\in[|\mathcal P|]$ such that $\{x_0\}\times[a]\subseteq P_q$.
\paragraph{Case 1: There is a good input.} Suppose that there exists a good input $x_0^\star\in[k]$ such that $\{x_0^\star\}\times[a]\subseteq P_q$ for some $q$. Let $\mathcal L_{x_0^\star}$ be the set of leaves of $\pi$ that contain an input where $x_0=x_0^\star$, that is,
\[\mathcal L_{x_0^\star}=\{\ell\in\mathcal L:\exists(x_1,y_0,y_1,z,i)\text{ with }\pi(x_0^*,x_1;y_0,y_1,z,i)=\ell\}\]
and let $\overline{\mathcal L_{x_0^\star}}$ denote the complementary set of leaves, that is $\overline{\mathcal L_{x_0^\star}}=\mathcal L\setminus\mathcal L_{x_0^\star}$. We will show that $|\mathcal L_{x_0^\star}|\ge L_d^B(f)$ and that $|\overline{\mathcal L_{x_0^\star}}|\ge2k-2$. As a result, we get that
\[L(\pi)\ge|\mathcal L_{x_0^\star}|+\overline{\mathcal L_{x_0^\star}}\ge L_d^B(f)+2k-2,\]
as desired.
First, we show that $|\mathcal L_{x_0^\star}|\ge L_d^B(f)$. Let $\pi'$ be the $d$-round Bob-first subprotocol of $\pi$ given when Alice says that her input is in $P_q$ at the first round. Since $\{x_0^\star\}\times[a]\subseteq P_q$, it follows that $\pi'$ computes $f(x_1,y_1)$ on input $(x_0^\star,x_1;y_0,y_1,0,1)$ for all $x_1\in[a]$ and $y_1\in[b]$. This yields a $d$-round Bob-first protocol for computing $f$, and therefore, the number of leaves in $\pi'$ containing the input $x_0^\star$ must be at least $L_d^B(f)$. Hence, $|\mathcal L_{x_0^\star}|\ge L_d^B(f)$.
Next, we argue that $|\overline{\mathcal L_{x_0^\star}}|\ge2k-2$. Consider the set of leaves given by
\[\{\pi(x_0,x_1;y_0,y_1,z,i):i=0,x_0=y_0\in[k]\setminus\{x_0^\star\},y_1=1,z\in\{0,1\}\}.\]
If we consider the restriction of $\pi$ to the inputs $y_1=1$ and $i=0$, we can apply Lemma~20 to conclude that all $2k-2$ leaves in this set are in $\overline{\mathcal L_{x_0^\star}}$ and are pairwise distinct.
\paragraph{Case 2: No good input.} Now we consider the case where there is no good input. For each $x_0\in[k]$, we define a set $A_{x_0}\subseteq[a]$ of cardinality 2 as follows. Since $x_0\in[k]$ is not good, there exists a pair $(x_1,x'_1)\in[a]^2$ such that $\pi(x_0,x_1;y_0,y_1,z,i)\ne\pi(x_0,x'_1;y_0,y_1,z,i)$ for all $y_0,y_1,z$ and $i$. Let $A_{x_0}=\{x_1,x'_1\}$. This completes our definition of $A_{x_0}$.
We claim that the $4k$ inputs in the following set are all in pairwise distinct leaves:
\[W=\{(x_0,x_1,y_0,y_1,z,i):x_0=y_0\in[k],i=0,x_1\in A_{x_0},y_1=1,z\in\{0,1\}\}.\]
To see this, suppose that $w\ne w'$ for some $w,w'\in W$. Let $w=(x_0,x_1;y_0,y_1,z,i)$ and $w'=(x'_0,x'_1;y'_0,y'_1,z',i')$. We prove $\pi(w)\ne\pi(w')$ by considering the following three cases.
\begin{enumerate}
\item If $x_0=x'_0$ and $x_1\ne x'_1$, then we know that $\{x_1,x'_1\}=A_{x_0}$. By the construction of $A_{x_0}$, we can conclude that $\pi(w)\ne\pi(w')$.
\item If $x_0=x'_0$ and $x_1=x'_1$, then we must have $x_0=y_0=x'_0$ and $z\ne z'$ since $w\ne w'$; using Lemma~20, we conclude that $\pi(w)\ne\pi(w')$.
\item Lastly, suppose that $x_0\ne x'_0$. If we consider the restriction of $\pi$ to the inputs $y_1=1$ and $i=0$, we can apply Lemma~20 to conclude that $\pi(w)\ne\pi(w')$.
\end{enumerate}
\end{proof}
Using this lower bound, we show one can approximately compute round-$d$ complexity given an oracle that approximately computes round-$(d+1)$ complexity. We consider the following notion of approximation.
\begin{definition}
For every constant $\epsilon>0$, we say that an oracle $\mathcal O$ is a $(1+\epsilon)$-approximation of a function $\mathcal L(\cdot)$ if there exists a constant $c$ such that, for all $g$ in the domain of $\mathcal L(\cdot)$,
\[\mathcal L(g)\le\mathcal O(g)\le(1+\epsilon)\cdot\mathcal L(g)+c.\]
\end{definition}
\begin{corollary}\label{cor:approx-elimination}
Let $0<\epsilon<\frac18$. Given an oracle computing a $(1+\epsilon)$-approximation of $L_{d+1}^A(\cdot)$ and a function $f:[a]\times[b]\to\{0,1\}$, one can deterministically compute a $(1+4\epsilon)$-approximation of $L_d^B(f)$ in time $(ab)^{O(1)}$.
\end{corollary}
\begin{proof}
First, we give the reduction algorithm and then we prove its correctness. Suppose $\mathcal O$ is an oracle that computes an approximation of $L_{d+1}^A(\cdot)$ satisifying, for all functions $g$,
\[L_{d+1}^A(g)\le\mathcal O(g)\le(1+\epsilon)L_{d+1}^A(g)+O(1).\]
For a positive integer $k$, let $F_k:([k]\times[a])\times([k]\times[b]\times\{0,1\}\times\{0,1\})\to\{0,1\}$ be given by
\[F_k(x_0,x_1;y_0,y_1,z,i)=\begin{cases}\XorEq_k(x_0;y_0,z)&\text{if }i=0,\\f(x_1;y_1)&\text{if }i=1.\end{cases}\]
The reduction computes
\[v:=\max\{\mathcal O(F_k)-2(1+\epsilon)k:k\in[ab]\},\]
and outputs $v':=(v+2)/(1-2\epsilon)$. It is easy to see that this reduction runs in time $(ab)^{O(1)}$.
To prove the correctness of the reduction, we claim that
\[L_d^B(f)\le v'\le(1+4\epsilon)\cdot L_d^B(f)+O(1)\]
for all functions $f$.
From Lemma~21, we know that for all $k$
\begin{align*}
\mathcal O(F_k)-2(1+\epsilon)k
&\le(1+\epsilon)L_{d+1}^A(F_k)-2(1+\epsilon)k+O(1)\\
&\le(1+\epsilon)L_d^B(f)+O(1)
\end{align*}
so we have that
\[v'=\frac{v+2}{1-2\epsilon}\le\frac{1+\epsilon}{1-2\epsilon}\cdot L_d^B(f)+O(1)\le(1+4\epsilon)\cdot L_d^B(f)+O(1),\]
where the last inequality holds because $\epsilon<1/8$.
On the other hand, if $k=L_d^B(f)$, we have from Lemma~21 that
\begin{align*}
\mathcal O(F_k)-2(1+\epsilon)k
&\ge L_{d+1}^A(F_k)-2(1+\epsilon)k\\
&\ge\min\{4k,2k-2+L_d^B(f)\}-2(1+\epsilon)k\\
&=\min\{4L_d^B(f),3L_d^B(f)-2\}-2(1+\epsilon)L_d^B(f)\\
&=3L_d^B(f)-2-2(1+\epsilon)L_d^B(f)\\
&\ge(1-2\epsilon)L_d^B(F_k)-2.
\end{align*}
Since $k=L_d^B(f)\le ab$, we conclude that $v'=(v+2)/(1-2\epsilon)\ge L_d^B(F_k)$.
\end{proof}
Combining Corollary~23 with the hardness result for the $d=3$ case in Theorem~18, we get that computing $L_d$ is NP-hard (under a polynomial-time truth-table reduction).
